% ============================================================
%  Yelp 餐饮评论精简深化分析报告
%  编译方式: xelatex -> bibtex -> xelatex -> xelatex
%  依赖宏包: ctex, geometry, graphicx, booktabs, longtable,
%            hyperref, amsmath, multirow, xcolor, listings,
%            caption, subcaption, float, fancyhdr, titlesec
% ============================================================
\documentclass[12pt, a4paper]{article}

% ── 中文支持 ─────────────────────────────────────────────────
\usepackage[UTF8, heading=true]{ctex}
\setCJKmainfont{SimSun}[BoldFont=SimHei, ItalicFont=KaiTi]
\setCJKsansfont{Microsoft YaHei}

% ── 页面布局 ─────────────────────────────────────────────────
\usepackage[
  top=2.5cm, bottom=2.5cm,
  left=3.0cm, right=2.5cm,
  headheight=15pt
]{geometry}

% ── 基础宏包 ─────────────────────────────────────────────────
\usepackage{amsmath, amssymb}
\usepackage{graphicx}
\usepackage{xcolor}
\usepackage{booktabs}
\usepackage{longtable}
\usepackage{multirow}
\usepackage{float}
\usepackage{caption}
\usepackage{subcaption}
\usepackage{enumitem}
\usepackage{fancyhdr}
\usepackage{titlesec}
\usepackage{tcolorbox}
\usepackage{array}
\usepackage{tabularx}
\usepackage{makecell}
\usepackage{diagbox}
\usepackage[numbers, sort&compress]{natbib}
\usepackage[
  colorlinks=true,
  linkcolor=NavyBlue,
  citecolor=NavyBlue,
  urlcolor=NavyBlue
]{hyperref}
\newcommand{\doi}[1]{\href{https://doi.org/#1}{\texttt{doi:#1}}}

% ── 自定义颜色 ───────────────────────────────────────────────
\definecolor{NavyBlue}{RGB}{25, 77, 140}
\definecolor{DarkGreen}{RGB}{39, 174, 96}
\definecolor{DarkRed}{RGB}{192, 57, 43}
\definecolor{LightGray}{RGB}{245, 246, 250}
\definecolor{AccentBlue}{RGB}{41, 128, 185}
\definecolor{AccentOrange}{RGB}{243, 156, 18}
\definecolor{BoxBorder}{RGB}{52, 152, 219}

% ── 页眉页脚 ─────────────────────────────────────────────────
\pagestyle{fancy}
\fancyhf{}
\fancyhead[L]{\small\color{gray} Yelp 餐饮评论精简深化分析报告}
\fancyhead[R]{\small\color{gray} \today}
\fancyfoot[C]{\thepage}
\renewcommand{\headrulewidth}{0.4pt}

% ── 节标题样式 ───────────────────────────────────────────────
\titleformat{\section}
  {\large\bfseries\color{NavyBlue}}
  {\thesection}{1em}{}[\titlerule]
\titleformat{\subsection}
  {\normalsize\bfseries\color{AccentBlue}}
  {\thesubsection}{1em}{}
\titleformat{\subsubsection}
  {\normalsize\bfseries}
  {\thesubsubsection}{1em}{}

% ── 结论框 ───────────────────────────────────────────────────
\newenvironment{insightbox}[1][关键结论]{%
  \begin{tcolorbox}[
    colback=LightGray, colframe=BoxBorder,
    title=\textbf{#1},
    fonttitle=\bfseries\color{white},
    coltitle=white,
    boxrule=0.8pt, arc=3pt,
    left=6pt, right=6pt, top=4pt, bottom=4pt
  ]
}{%
  \end{tcolorbox}
}

% ── 标题信息 ─────────────────────────────────────────────────
\title{
  \vspace{-1cm}
  {\huge\bfseries\color{NavyBlue}
    Yelp 餐饮评论数据的多维深度分析与运营决策研究}\\[0.5em]
  {\large
    基于 Wilson Score 排名、STL 时序分解与 LDA 主题建模的综合分析框架}
}
\author{
  \large 2023111348\quad 王岩方\\[0.3em]
  \small 数据来源：Yelp Open Dataset（Johnnyeee/Yelpdata\_663，HuggingFace）\\
  \small 数据规模：29,997 条评论 $\cdot$ 8,424 家餐厅 $\cdot$ 503 个城市 $\cdot$ 2004--2020\\[0.4em]
  \small 项目代码：\url{https://github.com/2023111348wangyf/restaurant_analysis}
}
\date{\today}

% ============================================================
\begin{document}
% ============================================================

\maketitle
\thispagestyle{empty}

% ── 摘要 ────────────────────────────────────────────────────
\begin{abstract}
\noindent
随着在线点评平台的普及，用户生成的餐饮评论已成为消费决策与运营改进的重要数据来源。
本研究以 Yelp 公开餐饮评论数据集为基础，
对 29,997 条英文评论、8,424 家餐厅及 503 个城市的用户生成内容进行系统性挖掘，
构建覆盖\textbf{表格}、\textbf{序列}与\textbf{文本}三类数据维度的综合分析框架。

在\textbf{表格维度}，采用 Wilson Score 置信下界（95\%）对餐厅进行统计意义明确的排名，
并运用 KMeans 算法将 8,424 家餐厅聚类为 Premium/Good/Average/At Risk 四个经营层级；
在\textbf{序列维度}，引入 STL 季节趋势分解（周期 $= 12$ 月）对月度评论量进行
趋势--季节--残差三分量解耦，结合线性回归外推短期活跃度预测，
并以数据驱动的 KMeans 聚类优化传统规则阈值 RFM 用户分层；
在\textbf{文本维度}，部署 LDA 主题建模（$k=8$）发现评论语料中的隐含话题结构，
训练 TF-IDF + 逻辑回归情感分类器（测试集准确率约 84\%），
并构建维度--情感矩阵量化 Food/Service/Ambiance/Price/Location
五维服务质量的情感格局。

研究结果表明：Wilson Score 在 Top 20--100 区间与贝叶斯平滑产生约 12\% 的位次差异；
服务维度（差评率 27\%）与等待时间话题（负面占比 52\%）
是最优先的运营改善切入点；KMeans RFM 相较规则分层可额外识别约 4.3\% 的高风险流失用户。
上述分析结果结合 P0--P3 分级差评预警体系，
可直接支撑平台推荐排序、精准用户运营及 NLP 模型微调数据集构建等实际应用场景。

\medskip
\noindent\textbf{关键词：}
餐饮评论挖掘；Wilson Score；贝叶斯排名；KMeans 聚类；
STL 趋势分解；时序预测；RFM 用户分层；
LDA 主题建模；TF-IDF；逻辑回归情感分类；ABSA；差评预警
\end{abstract}

\tableofcontents
\newpage

% ============================================================
\section{引言}
% ============================================================

\subsection{研究背景}

随着 Yelp、大众点评等在线点评平台的规模化发展，
用户生成内容（User-Generated Content, UGC）已成为餐饮消费决策的核心信息来源。
Yelp 作为北美最具影响力的餐饮评论平台，
积累了覆盖食物、服务、环境、价格等多个维度的亿级评论数据，
为餐饮企业的数字化运营提供了丰富的数据基础。

然而，海量评论数据的有效利用面临三重挑战：
其一，\textbf{评分可信度}——评论数量差异悬殊时，简单均分排名存在显著统计偏差；
其二，\textbf{时序结构}——评论量随平台发展与外部事件（如疫情）呈现复杂的趋势与季节效应，
简单聚合掩盖了关键动态信息；
其三，\textbf{语义洞察}——服务质量的短板与用户关注的隐含话题难以从结构化数据中直接提取。

本研究针对上述三类问题，
分别引入统计估计、时序分析与自然语言处理领域的成熟方法，
构建以\textbf{表格数据}、\textbf{序列数据}、\textbf{文本数据}为主轴的综合分析框架，
并将分析结果转化为可执行的运营决策建议。

\subsection{核心研究问题}

\begin{enumerate}[leftmargin=2em]
  \item \textbf{排名公平性（表格）}：Wilson Score 相对于贝叶斯平滑的改进体现在哪里？
        餐厅聚类能否发现隐藏的结构性分层？
  \item \textbf{时序规律（序列）}：STL 如何将长期趋势与季节效应解耦？
        KMeans RFM 与规则 RFM 的用户分布有何差异？
  \item \textbf{语义洞察（文本）}：LDA 发现哪些规律性话题？
        本地训练的情感分类器达到何种精度？
        各服务维度的情感结构有何差异？
\end{enumerate}

\subsection{数据集说明}

本研究使用的数据集来源于 Yelp 官方学术开放数据集 \cite{yelp_academic}，
经 Johnnyeee \cite{yelp_dataset} 整理为 HuggingFace Parquet 格式后采样获得。
经 \texttt{load\_yelp.py} 处理后，
保留 29,997 条有效英文评论、8,424 家餐厅、503 个城市，时间跨度 2004--2020 年。
情感标注规则：4--5 星→正面（+1），3 星→中性（0），1--2 星→负面（$-1$）。

% ============================================================
\section{数据预处理}
% ============================================================

\subsection{字段标准化与衍生特征}

原始 Yelp 数据经以下步骤统一为项目 schema：

\begin{enumerate}[leftmargin=2em]
  \item \textbf{字段映射}：\texttt{stars\_y} $\to$ \texttt{review\_score}，
        \texttt{text} $\to$ \texttt{review\_content}，
        \texttt{useful} $\to$ \texttt{like\_count}。
  \item \textbf{类型转换}：评分裁剪至 $[1,5]$，时间转为 \texttt{datetime}。
  \item \textbf{衍生特征}：评论文本长度（\texttt{content\_len}）、
        时间特征（年/月/小时/星期）、ABSA 维度数量（\texttt{aspect\_count}）。
  \item \textbf{异常过滤}：移除文本长度 $< 5$ 字符的无效评论，
        保留 \textbf{29,997} 条有效记录。
\end{enumerate}

\subsection{基础统计概览}

\begin{table}[H]
\centering
\caption{数据集基础统计（预处理后）}
\label{tab:basic_stats}
\begin{tabular}{llll}
\toprule
\textbf{指标} & \textbf{值} & \textbf{指标} & \textbf{值} \\
\midrule
总评论数   & 29,997 条  & 平均评论长度 & $\approx$ 685 字符 \\
餐厅数     & 8,424 家   & 评分均值     & 3.76 分 \\
用户数     & 27,672 人  & 评分标准差   & 1.38 分 \\
城市数     & 503 个     & 正面评论占比  & 68.1\% \\
州数       & 47 个      & 负面评论占比  & 20.3\% \\
时间跨度   & 2004--2020 & 中性评论占比  & 11.6\% \\
\bottomrule
\end{tabular}
\end{table}

评论评分分布呈明显的两极化特征：5 星评论约占 52\%，
1--2 星评论约占 20\%，中间评分（3 星）占比最低，
构成 Hu et al. \cite{hu2009overcoming} 所描述的"J 型分布"（J-shaped distribution）。
这种选择性偏差源于用户在极度满意或极度不满时的写评冲动显著高于中性体验，
在建模时需注意类别不平衡问题（见图 \ref{fig:rating_dist}）。

\begin{figure}[htbp]
\centering
\includegraphics[width=0.72\textwidth]{../outputs/figures/fig_rating_distribution.pdf}
\caption{评论评分分布直方图。评分呈右偏"J 型"，5 星占比最高（约 52\%），
         中性 3 星占比最低，正负样本比约 3.2:1。}
\label{fig:rating_dist}
\end{figure}

% ============================================================
\section{表格数据：排名深化与餐厅聚类}
% ============================================================

\subsection{Wilson Score 置信下界排名}

\subsubsection{方法原理与对比}

基于推荐系统中的评分可信度校正问题 \cite{bayesian_rating}，
贝叶斯平滑排名（式 \ref{eq:bayesian}）通过全局均值先验收缩样本均分，
但其惩罚强度由固定先验参数 $m$ 决定，缺乏概率区间解释。
Wilson Score \cite{wilson1927} 将均分归一化为二项比例，
直接给出在置信水平 $1-\alpha$ 下评分真值的统计下界：

\begin{equation}
\text{贝叶斯评分} = \frac{\mu \cdot m + \bar{r}_i \cdot n_i}{m + n_i}
\label{eq:bayesian}
\end{equation}

\begin{equation}
p = \frac{\bar{r}_i - 1}{4}, \quad
\text{Wilson 下界} =
\frac{p + \dfrac{z^2}{2n} - z\sqrt{\dfrac{p(1-p)}{n} + \dfrac{z^2}{4n^2}}}
     {1 + \dfrac{z^2}{n}}
\label{eq:wilson}
\end{equation}

其中 $z = 1.96$（$\alpha = 0.05$），计算结果映射回 $[1,5]$：
$\text{Wilson 评分} = \text{Wilson 下界} \times 4 + 1$。

\subsubsection{方法对比}

\begin{table}[H]
\centering
\caption{贝叶斯平滑 vs Wilson Score 方法对比}
\label{tab:ranking_compare}
\begin{tabularx}{\textwidth}{l X X}
\toprule
\textbf{维度} & \textbf{贝叶斯平滑} & \textbf{Wilson Score} \\
\midrule
统计基础 & 先验-后验均值加权，先验参数 $m=30$ 需手工设定 &
  二项比例置信下界，仅需置信水平 $\alpha$ \\
低评论量惩罚 & 强度固定（取决于 $m$），与 $n$ 非线性关系弱 &
  $n$ 越小区间越宽，惩罚自动增强，更符合统计直觉 \\
可解释性 & 先验参数意义直观但需领域经验 &
  "以 95\% 概率不低于此分"含义精确 \\
排名差异（rank\_drift） &
  \multicolumn{2}{c}{约 12\% 头部餐厅在两种方法下排名相差 $>$ 5 位} \\
\bottomrule
\end{tabularx}
\end{table}

\subsubsection{结果分析}

Wilson 排名对低评论量高分餐厅的惩罚更为显著：
仅有 10 条评论但均分 4.9 的餐厅，
其 Wilson 分通常比贝叶斯分低 0.15--0.25，
导致排名下滑 10--30 位。
两种排名在 Top 20 以内高度一致（重叠率约 85\%），
Top 20--100 区间差异明显——这正是平台排序策略最需精细化的区间。

\begin{figure}[htbp]
\centering
\includegraphics[width=\textwidth]{../outputs/figures/fig_ranking_comparison.pdf}
\caption{Wilson Score 与贝叶斯排名对比（Top 100 餐厅）。
         左图散点颜色深度代表排名漂移绝对值，对角线为零漂移基准；
         右图为漂移量频率分布，正值表示 Wilson 给出更低排名（更严格惩罚）。}
\label{fig:ranking_cmp}
\end{figure}

\begin{insightbox}
  \begin{itemize}[leftmargin=1.5em, itemsep=2pt]
    \item Wilson Score 提供统计意义明确的排序依据，适合平台搜索排序与 A/B 测试。
    \item 建议同时展示贝叶斯分与 Wilson 分，当两者差距 $> 0.3$ 时，
          向用户提示"评论数尚少，评分参考性有限"。
    \item rank\_drift $> 10$ 的餐厅（约占 8\%）是排名争议区，
          应优先补充人工审核。
  \end{itemize}
\end{insightbox}

\subsection{KMeans 餐厅四象限聚类}

\subsubsection{方法设计}

以每家餐厅的\textbf{均分}（avg\_score）、\textbf{评论量}（review\_count）、
\textbf{差评率}（neg\_rate）、\textbf{平均评论长度}（avg\_len）
为特征向量，经 \texttt{StandardScaler} 归一化后运行 KMeans（$k=4$，
$n\_\text{init}=10$，random\_state=42）\cite{sklearn}，
按质心均分降序将四个簇命名为：Premium / Good / Average / At Risk。

\subsubsection{聚类结果}

\begin{table}[H]
\centering
\caption{KMeans 餐厅聚类汇总（典型特征）}
\label{tab:cluster}
\begin{tabularx}{\textwidth}{l c c c c X}
\toprule
\textbf{层级} & \textbf{店铺数} & \textbf{均分} & \textbf{评论量} & \textbf{差评率} & \textbf{运营含义} \\
\midrule
Premium  & $\sim$820  & $\geq 4.3$ & 适中 & $\leq 8\%$  & 口碑标杆，重点维护与宣传 \\
Good     & $\sim$2100 & 3.9--4.3   & 高   & 10--18\%    & 主力市场，持续优化短板 \\
Average  & $\sim$3600 & 3.4--3.9   & 低   & 18--28\%    & 提升空间大，需聚焦改善 \\
At Risk  & $\sim$1900 & $\leq 3.4$ & 任意 & $\geq 28\%$ & 差评率高，需立即干预 \\
\bottomrule
\end{tabularx}
\end{table}

聚类揭示了一个重要发现：\textbf{高评论量不等于高质量}——
Good 簇的评论量中位数是 Premium 簇的 2.1 倍，
但均分低约 0.3 分，说明热门餐厅并非都是高口碑餐厅。
At Risk 簇的差评率中位数约为 Premium 簇的 3.5 倍，
差异远超随机波动水平（$p < 0.001$，Mann-Whitney U 检验）。

\begin{figure}[htbp]
\centering
\includegraphics[width=\textwidth]{../outputs/figures/fig_restaurant_clusters.pdf}
\caption{KMeans 餐厅聚类结果（$k=4$）。
         左图展示均分与评论量对数的联合分布，颜色代表聚类标签；
         右图为各聚类餐厅数量及对应均分，揭示"高评论量不等于高口碑"的结构性规律。}
\label{fig:clusters}
\end{figure}

\begin{insightbox}
  \begin{itemize}[leftmargin=1.5em, itemsep=2pt]
    \item Premium 层仅占全体约 10\%，是平台"精选榜单"的核心候选池。
    \item At Risk 层约占 23\%，差评率超阈值，建议优先触发 P1/P2 差评预警。
    \item 聚类标签可与 RFM 用户分层交叉，识别"高价值用户对低评级餐厅的负面评论"
          这类高风险信号。
  \end{itemize}
\end{insightbox}

% ============================================================
\section{序列数据：趋势分解与深化用户分层}
% ============================================================

\subsection{时段性基础规律}

评论量的小时分布呈双峰特征（图 \ref{fig:hourly}）：
午间峰（12:00--14:00）占全日约 15\%，
晚间主峰（18:00--22:00）占约 38\%。
月度数据显示 2014--2018 年为增长期，
2019 年后增速趋缓，2020 年受 COVID-19 影响骤降。

\begin{figure}[htbp]
\centering
\includegraphics[width=0.9\textwidth]{../outputs/figures/fig_hourly_pattern.pdf}
\caption{评论量小时分布柱状图。橙色区为午餐高峰（11--14 时），
         红色区为晚餐高峰（18--22 时），体现典型的餐饮消费"双峰"节律。}
\label{fig:hourly}
\end{figure}

\subsection{STL 趋势-季节-残差分解}

\subsubsection{方法原理}

Seasonal and Trend decomposition using Loess（STL）\cite{cleveland1990stl}
是时序分析的工业级标准方法，将观测序列分解为：

\begin{equation}
Y_t = T_t + S_t + R_t
\label{eq:stl}
\end{equation}

其中 $T_t$ 为趋势分量（LOESS 平滑），$S_t$ 为季节分量（周期 $=12$ 月），
$R_t$ 为残差。本研究采用 \texttt{robust=True} 配置以抑制异常点的影响。
数据长度不足 24 期时自动降级为 12 期中心移动平均（CMA）。

\subsubsection{分解结果分析}

STL 分解将月度评论量分解为：
\begin{itemize}[leftmargin=2em]
  \item \textbf{趋势分量}：清晰呈现 2004--2018 年的线性增长阶段
        与 2019--2020 年的平台期/衰退期，消除了季节波动对趋势判断的干扰。
  \item \textbf{季节分量}：夏季（6--8 月）普遍高于全年均值约 6--10\%，
        元旦前后出现明显低谷，与餐饮消费的季节性规律一致。
  \item \textbf{残差分量}：2020 年 3--6 月残差绝对值达历史最大，
        对应疫情封锁期的异常冲击，验证了模型的分离能力。
\end{itemize}

\begin{figure}[htbp]
\centering
\includegraphics[width=\textwidth]{../outputs/figures/fig_stl_decomposition.pdf}
\caption{月度评论量 STL 分解四面板图。从上至下依次为：原始观测值、趋势分量（LOESS 平滑）、
         季节分量（周期 12 月）、残差分量；红色阴影区域标注 2020 年 COVID-19 冲击。}
\label{fig:stl}
\end{figure}

\subsection{线性趋势外推预测}

在 STL 趋势分量上拟合线性回归（\texttt{LinearRegression}），
外推未来 6 个月评论量预期。
由于 2020 年受疫情影响趋势断裂，
实际应用时建议仅以 2017--2019 年趋势段为拟合窗口，
以排除结构性断点对预测的扭曲。

预测结果为运营决策提供量化基准：
当预测评论量持续低于趋势线 2 个标准差时，
应触发平台层面的复苏干预（如激励活动、广告投放）。

\subsection{KMeans 数据驱动 RFM 用户聚类}

RFM（Recency-Frequency-Monetary）框架是客户价值细分的经典方法 \cite{rfm_model}，
通过近度、频度、货币三个维度对用户进行差异化定位。
本研究将 Monetary 适配为历史评论平均星级，对传统框架进行场景化扩展。

\subsubsection{方法对比：规则阈值 vs KMeans}

\begin{table}[H]
\centering
\caption{规则 RFM 与 KMeans RFM 方法对比}
\label{tab:rfm_compare}
\begin{tabularx}{\textwidth}{l X X}
\toprule
\textbf{维度} & \textbf{规则四分位 RFM} & \textbf{KMeans RFM} \\
\midrule
分层逻辑 & 固定四分位阈值 + 人工规则树，需领域知识设计边界 &
  特征标准化后 $k$-means 自动发现高密度簇，边界由数据驱动 \\
层级数量 & 固定 5 层（Champions/Loyal/Recent/Potential/At Risk） &
  灵活设置（本研究 $k=5$），簇标签由质心综合分命名 \\
鲁棒性 & 四分位边界随数据分布移动，跨期可比性差 &
  Silhouette 系数可定量评估聚类质量，支持参数调优 \\
可解释性 & 规则透明，业务人员易理解 & 需结合质心坐标说明各簇特征 \\
\bottomrule
\end{tabularx}
\end{table}

\subsubsection{KMeans RFM 构建}

特征向量为 $[-R, F, M]$（$R$ 取负值使近期用户具有更大特征值），
经 \texttt{StandardScaler} 标准化后运行 KMeans（$k=5$）。
质心综合分 $= -R_{\text{centroid}} + F_{\text{centroid}} + M_{\text{centroid}}$
降序命名：Champions / Loyal / Potential / At Risk / Inactive。

\subsubsection{分层结果}

\begin{table}[H]
\centering
\caption{KMeans RFM 用户分层结果}
\label{tab:km_rfm}
\begin{tabularx}{\textwidth}{l c c c c X}
\toprule
\textbf{层级} & \textbf{用户数} & \textbf{占比} & \textbf{均 R（天）} & \textbf{均 F} & \textbf{运营建议} \\
\midrule
Champions & $\sim$2,900 & 10.5\% & $<200$  & $>4$ & 专属会员权益、探店合作 \\
Loyal     & $\sim$5,600 & 20.3\% & 200--500 & 3--4 & 积分激励、社区运营 \\
Potential & $\sim$6,800 & 24.6\% & 300--700 & 2--3 & 满减活动、内容触达 \\
At Risk   & $\sim$7,200 & 26.0\% & 600--1000 & 1--2 & 召回优惠、主动回访 \\
Inactive  & $\sim$5,172 & 18.7\% & $>1000$  & 1    & 低成本维持，择机唤醒 \\
\bottomrule
\end{tabularx}
\end{table}

对比规则 RFM 与 KMeans RFM，At Risk 层在 KMeans 中占比更高（26\% vs 24.4\%），
说明数据驱动方法识别出了更多规则方法遗漏的隐性流失风险用户。

\begin{insightbox}
  \begin{itemize}[leftmargin=1.5em, itemsep=2pt]
    \item STL 分解成功隔离了疫情冲击（2020 年残差异常），
          趋势分量可为平台长期运营规划提供无偏参考。
    \item KMeans RFM 将用户流失识别率从规则方法的 24.4\% 提升至约 26\%，
          可额外挽回约 1,100 名高风险用户。
    \item 建议每季度重新运行 KMeans 更新聚类中心，
          以适应平台活跃度的季节漂移。
  \end{itemize}
\end{insightbox}

\begin{figure}[htbp]
\centering
\includegraphics[width=\textwidth]{../outputs/figures/fig_rfm_clustering.pdf}
\caption{KMeans RFM 用户聚类结果（$k=5$）。
         左图以近度（R）与频度（F）为坐标，颜色区分五类用户；
         右图为各聚类用户规模及占比，Champions 与 Loyal 合计约 30\% 的核心用户贡献了平台绝大部分内容。}
\label{fig:rfm}
\end{figure}

% ============================================================
\section{文本数据：主题建模、情感分类与维度矩阵}
% ============================================================

\subsection{LDA 主题建模}

\subsubsection{方法原理}

隐含狄利克雷分布（LDA）\cite{blei2003lda}
假设每篇文档由若干隐含主题的混合生成，
每个主题对应词汇表上的多项分布。
本研究使用 scikit-learn \cite{sklearn} 提供的 \texttt{CountVectorizer}
（max\_features=3000，去除英文停用词）构建词袋矩阵，
在线学习 LDA（$n\_\text{topics}=8$，\texttt{max\_iter}=30）。

\subsubsection{主题发现}

LDA 从 29,997 条评论中自动抽取 8 个语义主题，
主要涵盖以下类型（基于 top-10 关键词人工标注）：

\begin{table}[H]
\centering
\caption{LDA 主题分布（典型关键词）}
\label{tab:lda}
\begin{tabular}{clc}
\toprule
\textbf{主题 ID} & \textbf{代表关键词} & \textbf{评论占比} \\
\midrule
T0 & food, delicious, fresh, chicken, rice, sauce, dish, flavor & $\sim$22\% \\
T1 & service, staff, waiter, rude, slow, friendly, manager      & $\sim$18\% \\
T2 & place, great, love, back, definitely, recommend, next time  & $\sim$16\% \\
T3 & price, expensive, worth, value, money, bill, affordable     & $\sim$12\% \\
T4 & wait, time, hour, minutes, long, seated, reservation       & $\sim$10\% \\
T5 & atmosphere, decor, cozy, noise, clean, interior, ambiance  & $\sim$9\%  \\
T6 & burger, pizza, sandwich, fries, cheese, bread, order       & $\sim$8\%  \\
T7 & sushi, roll, Japanese, Thai, Chinese, noodle, spicy        & $\sim$5\%  \\
\bottomrule
\end{tabular}
\end{table}

主题结构揭示了用户关注的层级：\textbf{食物本身}（T0，22\%）是第一核心，
\textbf{服务体验}（T1+T4，28\%，含等待时间）合计超过食物，
说明服务问题是评论写作的首要动机之一。
\textbf{菜系偏好}（T6+T7）虽占比较低，
但为餐厅品类定向运营提供了细分依据。

\subsubsection{主题-情感交叉分析}

将主导主题与情感标签交叉：T4（等待相关）的负面评论占比约 52\%，
远高于 T2（整体正面体验，负面仅约 8\%），
印证了等待时间管理是服务短板中最直接的改进突破口。

\begin{figure}[htbp]
\centering
\includegraphics[width=0.82\textwidth]{../outputs/figures/fig_lda_topics.pdf}
\caption{LDA 8 主题比例分布（水平条形图，按比例降序排列）。
         食物（T0）与整体体验（T2）合计占约 38\%，
         服务+等待（T1+T4）占 28\%，体现评论写作的核心关注点。}
\label{fig:lda}
\end{figure}

\subsection{TF-IDF + 逻辑回归情感分类器}

\subsubsection{模型设计}

\begin{table}[H]
\centering
\caption{情感分类器配置}
\label{tab:classifier}
\begin{tabular}{ll}
\toprule
\textbf{组件} & \textbf{配置} \\
\midrule
特征提取 & TF-IDF，1-2 gram，max\_features=8000，sublinear TF，min\_df=3 \\
分类器 & 逻辑回归，$C=1.0$，class\_weight=balanced，max\_iter=1000 \\
划分 & 训练/测试 = 8:2，随机种子 42，stratified 分层采样 \\
\bottomrule
\end{tabular}
\end{table}

采用 \texttt{class\_weight=balanced} 应对正负样本 3.2:1 的不平衡，
等价于对少数类（负面/中性）施加更大惩罚权重。

\subsubsection{评估结果}

\begin{table}[H]
\centering
\caption{情感分类器评估（测试集，约 5,996 条）}
\label{tab:clf_result}
\begin{tabular}{lcccc}
\toprule
\textbf{类别} & \textbf{Precision} & \textbf{Recall} & \textbf{F1} & \textbf{Support} \\
\midrule
Negative (--1) & $\sim$0.72 & $\sim$0.78 & $\sim$0.75 & $\sim$1,219 \\
Neutral (0)    & $\sim$0.48 & $\sim$0.43 & $\sim$0.45 & $\sim$696  \\
Positive (+1)  & $\sim$0.91 & $\sim$0.89 & $\sim$0.90 & $\sim$4,081 \\
\midrule
Macro avg      & $\sim$0.70 & $\sim$0.70 & $\sim$0.70 & --          \\
\textbf{Accuracy} & \multicolumn{4}{c}{$\sim$\textbf{0.84}} \\
\bottomrule
\end{tabular}
\end{table}

整体准确率约 84\%，正面类 F1 约 0.90，
中性类 F1 约 0.45（因中性评论与正/负界限模糊，此为预期偏低区间），
如图 \ref{fig:clf} 所示。相较于仅使用星级规则的方法（理论上限 100\%，但误标率约 8\%），
本分类器可实现对\textbf{无星级信号}的新评论进行实时情感打分，
具备实际部署价值。

\begin{figure}[htbp]
\centering
\includegraphics[width=0.82\textwidth]{../outputs/figures/fig_classifier_metrics.pdf}
\caption{情感分类器评估指标（TF-IDF + 逻辑回归，测试集 $\approx$ 5,996 条）。
         蓝色系柱图分别为 Precision / Recall / F1-Score，红色虚线为整体准确率（$\approx 0.84$）。
         中性类（Neutral）F1 最低，源于与正/负类的语义边界模糊。}
\label{fig:clf}
\end{figure}

\begin{insightbox}[情感分类器部署建议]
  \begin{itemize}[leftmargin=1.5em, itemsep=2pt]
    \item 将训练好的 \texttt{vectorizer + LogisticRegression} 模型序列化（pickle），
          接入评论入库流水线，对每条新评论实时打标。
    \item 中性类准确率偏低，建议将中性阈值改为
          $P(\text{neutral}) > 0.5$ 才归为中性，否则取 $\arg\max$ 的正/负类。
    \item 若需提升精度，可升级为 fine-tuned BERT \cite{devlin2019bert}（预期准确率 $\geq$91\%），
          本研究已按 7:1.5:1.5 比例导出三分训练集（共 29,955 条 JSONL），满足微调规模要求。
  \end{itemize}
\end{insightbox}

\subsection{维度-情感矩阵}

\subsubsection{方法}

对五大服务维度（Food / Service / Ambiance / Price / Location）
分别以关键词匹配提取命中评论子集，统计正/中/负情感比例，
构建 $5 \times 3$ 维度-情感矩阵。
该方法与 SemEval-2014 ABSA 任务 \cite{pontiki2014semeval}
的方面级情感分析范式一致，可直接为后续 ABSA 模型训练提供弱监督标注基础。

\begin{table}[H]
\centering
\caption{服务维度-情感矩阵（\%）}
\label{tab:asp_mat}
\begin{tabular}{lcccc}
\toprule
\textbf{维度} & \textbf{正面\%} & \textbf{中性\%} & \textbf{负面\%} & \textbf{总提及数} \\
\midrule
Food     & $\sim$76 & $\sim$6  & $\sim$18 & $\sim$24,900 \\
Service  & $\sim$61 & $\sim$12 & $\sim$27 & $\sim$20,100 \\
Ambiance & $\sim$66 & $\sim$12 & $\sim$22 & $\sim$12,600 \\
Price    & $\sim$58 & $\sim$18 & $\sim$24 & $\sim$11,400 \\
Location & $\sim$70 & $\sim$15 & $\sim$15 & $\sim$6,300  \\
\bottomrule
\end{tabular}
\end{table}

\subsubsection{结果解读}

\begin{itemize}[leftmargin=2em]
  \item \textbf{Food} 维度正面率最高（76\%），是品牌口碑的最强支撑。
  \item \textbf{Service} 维度差评率最高（27\%），服务标准化培训的口碑改善 ROI 最大。
  \item \textbf{Price} 维度正面率最低（58\%），但中性比例最高（18\%），
        说明价格感知存在较大主观方差，定价策略应配合明确的价值传递。
  \item \textbf{Location} 维度差评率最低（15\%），
        地理因素对用户体验的负面影响相对可控。
\end{itemize}

矩阵与 LDA 结果形成互补印证（见图 \ref{fig:asp_mat}）：
T1（服务）+T4（等待）共同解释了 Service 维度高差评率，
T3（价格）与 Price 维度的中性偏高相互验证。

\begin{figure}[htbp]
\centering
\includegraphics[width=0.68\textwidth]{../outputs/figures/fig_aspect_sentiment.pdf}
\caption{服务维度-情感矩阵热力图（$5 \times 3$，颜色越绿正面比例越高，越红负面比例越高）。
         Service 维度负面占比最高（约 27\%），Food 维度正面率最高（约 76\%），
         与 LDA 主题分析的结论相互印证。}
\label{fig:asp_mat}
\end{figure}

\begin{insightbox}[文本分析关键结论]
  \begin{itemize}[leftmargin=1.5em, itemsep=2pt]
    \item LDA 发现等待时间（T4）负面占比 52\%，是单一可改进因素中差评率最高的来源。
    \item TF-IDF + LR 分类器准确率约 84\%，可直接部署于评论实时打标流水线。
    \item 维度-情感矩阵将服务（27\%负面）定位为最高优先级改善方向，
          与 ABSA 数据集中服务类负面样本占比（30.3\%）高度一致。
    \item 后续建议：以 BERT 替代 TF-IDF+LR，
          训练集（29,955 条）已满足 fine-tuning 规模要求。
  \end{itemize}
\end{insightbox}

% ============================================================
\section{运营策略与行动建议}
% ============================================================

\subsection{差评预警体系}

四级滑动窗口预警体系（P0--P3）贯通三个分析维度：

\begin{table}[H]
\centering
\caption{差评预警分级体系}
\label{tab:alert}
\begin{tabularx}{\textwidth}{l c c X}
\toprule
\textbf{级别} & \textbf{窗口} & \textbf{差评率阈值} & \textbf{响应措施} \\
\midrule
P0（危机） & 1 天  & $\geq 50\%$ & 立即召开紧急会议，主动联系用户补偿 \\
P1（预警） & 3 天  & $\geq 35\%$ & 24h 内回复全部差评，复盘近期客诉 \\
P2（关注） & 7 天  & $\geq 25\%$ & 本周完成原因分析，推出口碑修复活动 \\
P3（跟踪） & 14 天 & $\geq 20\%$ & 每日监控趋势，与 KMeans 聚类结果对照 \\
\bottomrule
\end{tabularx}
\end{table}

\subsection{跨模块联动行动矩阵}

\begin{table}[H]
\centering
\caption{数据驱动运营行动建议（精简深化版）}
\label{tab:action}
\begin{tabularx}{\textwidth}{l l X}
\toprule
\textbf{分析模块} & \textbf{核心发现} & \textbf{行动建议} \\
\midrule
Wilson 排名  & rank\_drift $> 10$ 占 8\%   & 差异区间优先触发人工审核，向用户展示"评论数不足"提示 \\
KMeans 聚类  & At Risk 层占 23\%             & 交叉 P1/P2 预警，差评率触发自动生成整改工单 \\
STL 分解     & 疫情残差异常可隔离            & 排除 2020 年异常后外推趋势，作为运营目标基线 \\
KMeans RFM   & At Risk 用户多识别 1,100 名   & 30 天内定向召回，目标回访率 $\geq 15\%$ \\
LDA T4 主题  & 等待相关负面占 52\%           & 等待超 20 分钟自动触发安抚短信 + 优惠券 \\
情感分类器   & 准确率 $\sim$84\%             & 接入评论入库流水线，打标结果写入 sentiment\_score 字段 \\
维度矩阵     & Service 差评率最高（27\%）    & 制定服务 SOP，服务口碑纳入员工 KPI 考核 \\
ABSA 数据集  & 55,821 条细粒度标注           & 微调 BERT 模型，替换 LR 分类器，提升实时识别精度 \\
\bottomrule
\end{tabularx}
\end{table}

% ============================================================
\section{结论}
% ============================================================

本报告对 Yelp 餐饮评论数据实施\textbf{精简深化}重构，
在移除信息冗余的图数据与时空数据模块后，
对表格、序列、文本三个核心维度引入了更严格的统计与机器学习方法，
得出以下核心结论：

\begin{enumerate}[leftmargin=2em]
  \item \textbf{排名公平性}：Wilson Score 为低评论量店铺提供了统计意义明确的惩罚机制，
        其排名结果在 Top 20--100 区间与贝叶斯平滑产生约 12\% 的显著差异，
        更适合作为平台推荐排序的生产依据。

  \item \textbf{餐厅分层}：KMeans 聚类揭示高评论量不等于高口碑的结构性规律，
        At Risk 层（约 23\%）的差评率显著高于其他聚类（$p < 0.001$），
        是差评预警体系的最高优先覆盖目标。

  \item \textbf{趋势解耦}：STL 分解成功分离了 2020 年疫情冲击（残差异常），
        季节分量显示夏季评论量高出均值约 6--10\%，
        趋势外推为平台运营提供了去季节化的量化基准。

  \item \textbf{用户分层精化}：KMeans RFM 相较于规则阈值多识别约 1,100 名
        At Risk 用户（+4.3\%），数据驱动分层在挽救流失用户方面具有实际价值。

  \item \textbf{主题挖掘}：LDA 发现等待时间（T4）主题负面占比高达 52\%，
        是单一可干预因素中差评率最高的来源，
        优先于服务态度（T1，负面约 30\%）成为最高价值改善点。

  \item \textbf{情感分类}：本地训练的 TF-IDF + 逻辑回归分类器准确率约 84\%，
        可直接部署于评论打标流水线；
        若精度要求更高，BERT 训练集（29,955 条）已就绪。
\end{enumerate}

\subsection{研究局限性}

\begin{itemize}[leftmargin=2em]
  \item 数据截止 2020 年，后疫情消费行为变化未能覆盖。
  \item 样本为 30,000 条随机采样，可能存在城市分布不均的抽样偏差。
  \item 时空维度的移除意味着地理竞争密度、租金等选址关键变量未被纳入，
        如需选址决策仍需补充外部数据源。
  \item LDA 主题的语义标注依赖人工解读，存在一定主观性。
\end{itemize}

% ============================================================
% 参考材料
% ============================================================
\newpage
\renewcommand{\refname}{参考材料}
\begin{thebibliography}{99}

% ── 数据来源 ─────────────────────────────────────────────────

\bibitem{yelp_academic}
Yelp Inc.
\textit{Yelp Open Dataset}.
Retrieved from \url{https://www.yelp.com/dataset}, accessed 2024.

\bibitem{yelp_dataset}
Johnnyeee.
\textit{Yelpdata\_663: Yelp Restaurant Review Dataset} [Data set].
Hugging Face, 2022.
\url{https://huggingface.co/datasets/Johnnyeee/Yelpdata_663}

% ── 统计与排名方法 ────────────────────────────────────────────

\bibitem{wilson1927}
Wilson, E. B. (1927).
Probable inference, the law of succession, and statistical inference.
\textit{Journal of the American Statistical Association}, \textbf{22}(158), 209--212.
\doi{10.2307/2276774}

\bibitem{bayesian_rating}
Resnick, P., \& Varian, H. R. (1997).
Recommender systems.
\textit{Communications of the ACM}, \textbf{40}(3), 56--58.
\doi{10.1145/245108.245121}

% ── 时间序列分析 ──────────────────────────────────────────────

\bibitem{cleveland1990stl}
Cleveland, R. B., Cleveland, W. S., McRae, J. E., \& Terpenning, I. (1990).
STL: A seasonal-trend decomposition procedure based on Loess.
\textit{Journal of Official Statistics}, \textbf{6}(1), 3--73.

% ── 用户分层 ─────────────────────────────────────────────────

\bibitem{rfm_model}
Fader, P. S., Hardie, B. G. S., \& Lee, K. L. (2005).
RFM and CLV: Using iso-value curves for customer base analysis.
\textit{Journal of Marketing Research}, \textbf{42}(4), 415--430.
\doi{10.1509/jmkr.2005.42.4.415}

% ── 自然语言处理 ──────────────────────────────────────────────

\bibitem{blei2003lda}
Blei, D. M., Ng, A. Y., \& Jordan, M. I. (2003).
Latent Dirichlet Allocation.
\textit{Journal of Machine Learning Research}, \textbf{3}, 993--1022.

\bibitem{devlin2019bert}
Devlin, J., Chang, M. W., Lee, K., \& Toutanova, K. (2019).
BERT: Pre-training of deep bidirectional transformers for language understanding.
In \textit{Proceedings of the 2019 Conference of the North American Chapter
of the Association for Computational Linguistics: Human Language Technologies
(NAACL-HLT 2019)}, Vol.\ 1, pp. 4171--4186.
Minneapolis, MN: Association for Computational Linguistics.
\doi{10.18653/v1/N19-1423}

\bibitem{pontiki2014semeval}
Pontiki, M., Galanis, D., Pavlopoulos, J., Papageorgiou, H.,
Androutsopoulos, I., \& Manandhar, S. (2014).
SemEval-2014 Task 4: Aspect based sentiment analysis.
In \textit{Proceedings of the 8th International Workshop on Semantic Evaluation
(SemEval 2014)}, pp. 27--35.
Dublin, Ireland: Association for Computational Linguistics.
\doi{10.3115/v1/S14-2004}

% ── 机器学习工具库 ────────────────────────────────────────────

\bibitem{sklearn}
Pedregosa, F., Varoquaux, G., Gramfort, A., Michel, V., Thirion, B.,
Grisel, O., Blondel, M., Prettenhofer, P., Weiss, R., Dubourg, V.,
Vanderplas, J., Passos, A., Cournapeau, D., Brucher, M., Perrot, M.,
\& Duchesnay, \'{E}. (2011).
Scikit-learn: Machine learning in Python.
\textit{Journal of Machine Learning Research}, \textbf{12}, 2825--2830.

% ── 评论偏差研究 ──────────────────────────────────────────────

\bibitem{hu2009overcoming}
Hu, N., Pavlou, P. A., \& Zhang, J. (2009).
Overcoming the J-shaped distribution of product reviews.
\textit{Communications of the ACM}, \textbf{52}(10), 144--147.
\doi{10.1145/1562764.1562800}

\end{thebibliography}

\end{document}
